\section{Validation}
\label{sec:validation}

The implementation runs two gated checks against the Fed's own artefacts, and they are evidence of very different strength. One is a \emph{software invariant} that is true by construction and carries no economic information; the other is a genuine cross-implementation comparison. They are presented in that order, and labelled as such, because the first was previously reported as this paper's headline validation result and it cannot bear that weight. All experiments use the VAR model over 2026Q1--2030Q4 with the Fed's standard demo fiscal configuration (surplus-ratio targeting: \texttt{dfpdbt} $=0$, \texttt{dfpsrp} $=1$) and the April 2026 LONGBASE vintage.

\subsection{The tracking invariant is an identity, not evidence}
\label{subsec:tracking}

After \texttt{init\_trac}, solving the baseline reproduces LONGBASE for all 284 endogenous variables in all 20 quarters to a maximum absolute error of $5.6\times10^{-17}$. \textbf{This number is true by construction and must not be quoted as evidence that the implementation agrees with the Fed's data.} Three facts establish why.

First, \texttt{init\_trac} sets each equation's add factor to \emph{minus that equation's residual evaluated at the input data}, so after tracking, \texttt{solve} is algebraically the identity map on whatever it was tracked to --- for any input whatsoever. Second, \texttt{solve\_periods} seeds the Newton iteration with the current row of the input data, i.e.\ with the answer: at the tracked baseline, 282 of the 284 equations have a \emph{bit-exactly zero} residual before the solver runs at all (the largest is $8.7\times10^{-19}$), so Newton converges without moving. Third, the residual that remains is floating-point non-associativity in re-evaluating $A - (A-B) - B$, not solver accuracy: 5{,}645 of the 5{,}680 compared cells are bit-identical.

The decisive demonstration is itself a gated test (\texttt{tests/test\_tracking.py::\allowbreak test\_\allowbreak tracking\_\allowbreak invariant\_\allowbreak holds\_\allowbreak for\_\allowbreak arbitrary\_\allowbreak data}): multiplying every endogenous series by an independent random factor drawn from $U[0.5,1.5]$ --- which destroys every accounting identity in the database ($C+I+G+NX$ no longer sums to GDP) and every estimated behavioural relation in the model --- still passes the identical gate. Table~\ref{tab:tracking} reports the gate on inputs of decreasing economic meaning; the last row is the point.

\begin{table}[htbp]
\centering
\caption{The tracking gate on inputs of decreasing economic meaning (2026Q1--2030Q4). An identity, not a measurement: the gate passes on data in which every accounting identity has been destroyed.}
\label{tab:tracking}
\begin{tabular}{lccc}
\toprule
Baseline handed to \texttt{init\_trac} & Max abs.\ error & Max rel.\ error & Passes the gate? \\
\midrule
LONGBASE (the published figure) & $5.6\times10^{-17}$ & $5.6\times10^{-17}$ & yes \\
LONGBASE $+$ 5\% i.i.d.\ noise on every cell & $4.4\times10^{-11}$ & $1.8\times10^{-13}$ & yes \\
Every endogenous series $\times\,U[0.5,1.5]$ & $6.7\times10^{-9}$ & $9.1\times10^{-12}$ & yes \\
\midrule
\multicolumn{2}{l}{\emph{Test gate}} & \multicolumn{2}{r}{$<10^{-8}$ abs., $<10^{-10}$ rel.} \\
\bottomrule
\end{tabular}

\medskip
\parbox{0.95\textwidth}{\footnotesize The test is worth keeping: it pins that \texttt{init\_trac} and \texttt{solve} remain exact inverses, which would catch a real regression in the add-factor machinery. It says nothing about whether the model's economics are right, and nothing about this implementation's numerical quality relative to any other. The substantive fidelity evidence in this paper is the cross-implementation comparison of Section~\ref{subsec:crossval}, where the add factors are frozen and the solver has genuine work to do, and the multiplier comparisons of Section~\ref{sec:multipliers}.}
\end{table}

\subsection{Cross-validation against the Fed's pyfrbus}
\label{subsec:crossval}

This is the check that carries the fidelity claim. Unlike the tracking invariant, it compares two independent codebases on a \emph{shocked} path, where the add factors are held fixed and both solvers must actually iterate to a fixed point; there is no construction that forces the two answers to agree.

The identical experiment --- \texttt{init\_trac} followed by a 100-basis-point shock to \texttt{rffintay\_aerr} in 2026Q1, mirroring the Fed's \texttt{demos/example1.py} --- was run in both this implementation and the Fed's \texttt{pyfrbus}, and all 284 endogenous variables were compared over all 20 quarters. The vendor reference was produced by \texttt{pyfrbus} running in a separate Python 3.11 environment with its Newton solver at \texttt{xtol} $=10^{-8}$; a documented script regenerates it, recording the out-of-repo patches needed to run the vendor package on modern dependencies (an \texttt{splu} shim replacing \texttt{scikit-umfpack}, unpinning \texttt{sympy}, a fix to the vendor's \texttt{Max}/\texttt{Min} derivative fallback under modern \texttt{symengine}, and \texttt{pandas} $<2$).

\begin{table}[htbp]
\centering
\caption{Cross-validation vs.\ pyfrbus: identical 100bp funds-rate shock, 2026Q1--2030Q4}
\label{tab:crossval}
\begin{tabular}{lcc}
\toprule
Metric & Value & Test gate \\
\midrule
Max absolute difference, all 284 endogenous $\times$ 20 quarters & $6.0\times10^{-9}$ & $<10^{-6}$ \\
Max relative difference (denominator clipped at $10^{-8}$) & $4.9\times10^{-8}$ & $<10^{-6}$ \\
\bottomrule
\end{tabular}
\end{table}

\begin{table}[htbp]
\centering
\caption{Cross-validation: agreement on key variables (max over 20 quarters)}
\label{tab:crossval-key}
\begin{tabular}{llcc}
\toprule
Variable & Description & Max abs.\ diff. & Max rel.\ diff. \\
\midrule
\texttt{xgdp}   & Real GDP                & $1.5\times10^{-10}$ & $5.6\times10^{-15}$ \\
\texttt{pcpi}   & CPI price level          & $8.5\times10^{-13}$ & $2.5\times10^{-15}$ \\
\texttt{picxfe} & Core PCE inflation       & $2.3\times10^{-13}$ & $1.2\times10^{-13}$ \\
\texttt{lur}    & Unemployment rate        & $4.0\times10^{-13}$ & $9.2\times10^{-14}$ \\
\texttt{rff}    & Federal funds rate       & $3.0\times10^{-13}$ & $1.1\times10^{-13}$ \\
\bottomrule
\end{tabular}
\end{table}

Tables~\ref{tab:crossval} and \ref{tab:crossval-key} summarise the agreement. The worst-agreeing variables are near-zero expectational gap terms (\texttt{zgap05}, relative difference $4.9\times10^{-8}$ on values of order $10^{-4}$), which reflect the two solvers' respective step tolerances rather than any semantic difference; with the vendor solver at its default looser tolerance (\texttt{scipy.optimize.root}, \texttt{xtol} $=10^{-6}$) the discrepancy is $\sim1.7\times10^{-5}$ and again concentrated in the same terms. The headline variables agree to $10^{-10}$ or better in absolute terms. Since the vendor's block-decomposed solver and this implementation's single-block Newton solve the same $F(x)=0$, agreement up to solver tolerance is the expected --- and observed --- outcome.

\subsection{How large is $10^{-8}$? The reference implementation against itself}
\label{subsec:refnoise}

A residual is only meaningful against a yardstick, and the natural one is the reference implementation's own version-to-version variation. The vendored reference is \texttt{pyfrbus} 1.0.0 (files dated 6~April~2022); the Board also ships 1.1.1 (5~February~2026), whose \texttt{model.xml} and LONGBASE are byte-identical to 1.0.0's, so a comparison between them isolates solver differences from model differences. Version 1.1.1 rewrites the Newton routine to reuse its LU factorisation across iterations, recomputing only when the residual fails to halve --- a change to the iterate path, not to the fixed point. Repeating the identical shock experiment across all three implementations gives Table~\ref{tab:refnoise}.

\begin{table}[htbp]
\centering
\caption{Three-way comparison: identical 100bp funds-rate shock, 284 endogenous $\times$ 20 quarters}
\label{tab:refnoise}
\begin{tabular}{lcc}
\toprule
Comparison & Max abs.\ diff. & Max rel.\ diff. \\
\midrule
This implementation vs.\ \texttt{pyfrbus} 1.0.0 & $6.0\times10^{-9}$ & $4.9\times10^{-8}$ \\
This implementation vs.\ \texttt{pyfrbus} 1.1.1 & $1.4\times10^{-8}$ & $1.8\times10^{-7}$ \\
\texttt{pyfrbus} 1.1.1 vs.\ \texttt{pyfrbus} 1.0.0 & $1.3\times10^{-8}$ & $1.3\times10^{-7}$ \\
\bottomrule
\end{tabular}
\end{table}

The third row is the informative one: the Board's own two releases differ from each other by as much as this implementation differs from either, and all three residuals are concentrated in the same near-zero series (\texttt{wpsn}, \texttt{zgap05}). The agreement documented above is therefore not an artefact of the particular vendored release --- it sits at the scale of the reference implementation's own numerical noise.

On the tracking invariant of Section~\ref{subsec:tracking}, \texttt{pyfrbus} 1.1.1 reproduces LONGBASE to $1.1\times10^{-8}$ where this implementation reports $5.6\times10^{-17}$. \textbf{That is not a quality ranking and an earlier version of this paper was wrong to present it as one.} Both figures are add-factor identities (Section~\ref{subsec:tracking}); the entire difference is that this implementation seeds Newton with the data row it is about to reproduce and hands it back unchanged, whereas the vendor's block-decomposed solver iterates each block to its own step tolerance. A smaller number here means less arithmetic was performed, not that the economics are closer to the Board's. The committed reference remains 1.0.0 so that the regression gate has a fixed anchor. Only the VAR path was exercised in this comparison; 1.1.1 additionally fixes a lead-substitution guard that affects the MCE path, which this implementation does not support and which was therefore not tested.
